Polarita molekuly závisí na \textbf{vektorovém součtu dipólových momentů všech vazeb}:
\[\boldsymbol p_{\mathrm{mol}}=\sum_i\boldsymbol p_i.\]
V SF$_6$ jsou vazby S--F polární, ale šest atomů F tvoří pravidelný oktaedr. Protilehlé vazby mají stejně velké a opačné dipólové momenty, které se vyruší. Těžiště kladného a záporného náboje splývají, proto molekula nemá permanentní dipólový moment a je navenek nepolární.
